\documentclass[../../../include/open-logic-section]{subfiles}

\begin{document}

\olfileid{sth}{ordinals}{vn}	
\olsection[Pambangunan von Neumann]{Pambangunan Ordinal Miturut von Neumann}

\olref[sth][ordinals][iso]{thm:woalwayscomparable}
ndadekake kita mikir. Kita bisa ngenalake obyek tartamtu,
kang diarani \emph{tipe urutan}, minangka wakil
kanggo urutan becik. Yen $\ordtype{A, <}$ ditulis
kanggo tipe urutan saka urutan becik $\tuple{A, <}$,
kita ngarepake rong prinsip iki:
\begin{align*}
	\ordtype{A, <} = \ordtype{B, \lessdot} & 
	\text{ yen lan mung yen } \ordeq{\tuple{A, <}}{\tuple{B, \lessdot}}\\
	\ordtype{A, <} < \ordtype{B, \lessdot}&
	\text{ yen lan mung yen }\ordeq{\tuple{A, <}}{\tuple{B_b, \lessdot_b}}\text{ kanggo sawenehing }b \in B
\end{align*}
Kajaba iku, kita bisa ngenalake tipe urutan
\emph{minangka himpunan tartamtu}, kaya anggone kita
ngenalake wilangan asli minangka himpunan tartamtu.

Cara kang paling lumrah kanggo nindakake iki---lan
iki kang bakal kita tindakake---yaiku netepake tipe
urutan kasebut liwat sawenehing himpunan kang diurutake
kanthi becik lan \emph{kanonik}. Himpunan kanonik iki
pisanan dikenalake dening von Neumann:

\begin{defn}
Himpunan $A$ iku \emph{transitif} yen lan mung yen
$(\forall x \in A)x \subseteq A$. Banjur $A$ iku
\emph{ordinal} yen lan mung yen $A$ transitif lan
diurutake kanthi becik dening $\in$.
\end{defn}
\noindent
Ing salajengipun, kita bakal nganggo aksara Yunani
kanggo ordinal. Saka definisi langsung ndherek yen
$\alpha$ ordinal, mula $\tuple{\alpha, \in_\alpha}$
iku urutan becik, kanthi $\in_\alpha =
\Setabs{\tuple{x, y} \in \alpha^2}{x \in y}$.
Dadi, kanthi rada nglonggari notasi, kita bisa kandha
yen $\alpha$ \emph{dhewe} iku urutan becik.

Iki sawetara ordinal kang kapisan:
\[
	\emptyset, \{\emptyset\}, 
	\{\emptyset, \{\emptyset\}\}, 
	\{\emptyset, \{\emptyset\}, \{\emptyset, \{\emptyset\}\}\}, \ldots
\]
Panjenengan bisa ndeleng yen iki padha karo ordinal
kapisan kang ditemoni ing Aksioma Tanpa Wates, yaiku
ing definisi $\omega$ miturut von Neumann (delengen
\olref[sth][z][infinity-again]{sec}). Iki dudu
bab kang dumadi tanpa sengaja. Definisi ordinal miturut von Neumann
nganggep wilangan asli minangka ordinal, nanging uga
ngidini ordinal transfinit.

Kaya biasane, saiki kita bisa takon: apa himpunan iki
\emph{pancen} ordinal? Utawa von Neumann mung menehi
sawenehing himpunan kang bisa kita \emph{anggep}
minangka ordinal? Rembugan bab pitakonan iki padha
jinise karo rembugan ing \olref[sfr][rel][ref]{sec},
\olref[sfr][arith][ref]{sec},
\olref[sfr][infinite][dedekindsproof]{sec}, lan
\olref[sth][z][nat]{sec}, mula kita ora bakal
ngrembug dawa ing kene. Ing salajengipun, kita bakal
nganggo ``ordinal'' kanggo nyebut ``ordinal von Neumann''.

\end{document}
